% !TEX program = pdflatex
\documentclass[3p,11pt]{elsarticle}
\usepackage{amsmath,amssymb}
\usepackage{array}
\usepackage{booktabs}
\usepackage[tableposition=top]{caption}
\usepackage{enumitem}
\usepackage{tikz}
\usepackage{pgfplots}
\pgfplotsset{compat=1.18}
\usepackage{url}
\urlstyle{same}
\usepackage[hidelinks]{hyperref}
\emergencystretch=3em
% first-page footer: the page number alone
\makeatletter
\def\ps@pprintTitle{\let\@oddhead\@empty\let\@evenhead\@empty
  \def\@oddfoot{\hfil\thepage\hfil}\let\@evenfoot\@oddfoot}
\makeatother

% ---- notation ---------------------------------------------------------------
\newcommand{\um}{\,$\mu$m}
\newcommand{\kms}{km\,s$^{-1}$}
\newcommand{\flux}{erg\,cm$^{-2}$\,s$^{-1}$}
\newcommand{\mA}{m\AA}
\newcommand{\sci}[2]{#1\times10^{#2}}
\newcommand{\ion}[2]{#1\,\textsc{#2}}
\newcommand{\nuc}[2]{$^{#1}$#2}
\newcommand{\lxl}{$L_X/L_\mathrm{bol}$}

\begin{document}

\begin{frontmatter}

\title{A search for laser lines in five large spectroscopic surveys of nearby stars}

\author{Charles C. Norton}

\begin{abstract}
A laser aimed at Earth from another star adds an unresolved emission line to the
star's spectrum. The large spectroscopic surveys now hold spectra of tens of thousands of
stars within $100$\,pc, taken for stellar parameters and abundances, and none has been
searched for such lines. We search five: the Gaia DR3 mean RVS spectra of $2363$ FGK
stars within $75$\,pc ($846$ to $870$\,nm), GALAH DR4 spectra of $621$ dwarfs within
$100$\,pc ($471$ to $789$\,nm in four windows), APOGEE DR17 combined spectra and
$14178$ visits of $3370$ dwarfs within $100$\,pc ($1.51$ to $1.70$\um, including the
$1530$ to $1565$\,nm band of terrestrial fibre-optic lasers), DESI DR1 spectra of $7705$
stars within $100$\,pc ($360$ to $982$\,nm), and LAMOST DR10 medium-resolution spectra
of $14400$ stars within $100$\,pc in $35179$ observations. Each spectrum is modelled by a
non-negative combination of the most similar spectra of its survey, so a narrow emission
line that they lack cannot be fitted away, and emission-like residuals are judged against
the absorption-like residuals of the same spectra. No emission line above that null
remains in the Gaia, APOGEE and multi-exposure DESI spectra; every GALAH candidate is
instrumental; and in LAMOST the $19$ features that recur in two exposures of single
stars are one or two pixels wide where the instrument's profile spans six, detector
pixels and not light. From each spectrum's noise and the star's magnitude, a line at the
detection threshold carries a median of $2\times10^{-16}$\,erg\,cm$^{-2}$\,s$^{-1}$ in
APOGEE and DESI and $1$ to $3\times10^{-15}$ in GALAH, Gaia and LAMOST, the flux of a
transmitter of $10^{19}$ to $10^{20}$\,W radiating isotropically or of $5$\,kW to
$70$\,kW behind a $10$\,m aperture at the median star, and of $5$ to $50$\,W at the
nearest.
\end{abstract}

\begin{keyword}
Technosignatures \sep Lasers \sep Techniques: spectroscopic \sep Solar neighbourhood \sep
Surveys
\end{keyword}

\end{frontmatter}

% =============================================================================
\section{Introduction}\label{sec:intro}
% =============================================================================

Schwartz and Townes proposed that a civilization could signal with optical masers, whose
light, confined to a narrow beam and a narrow band, would outshine its star at that
wavelength~\cite{SchwartzTownes1961}. A laser aimed at Earth would add an unresolved
emission line to the star's spectrum; from a $10$\,m aperture the beam is a few
hundredths of an astronomical unit wide at $10$\,pc and about one at $100$\,pc, so it
has to be held on the Earth's position when the light arrives. Searches have looked for
such lines in the high-resolution spectra of radial-velocity programmes: $577$ FGKM
stars with Keck/HIRES~\cite{ReinesMarcy2002}, $2796$ stars at $3640$ to
$7890$\,\AA~\cite{TellisMarcy2015}, $5600$ FGKM stars with thresholds of
megawatts~\cite{TellisMarcy2017}, $388$ stars with the Automated Planet
Finder~\cite{Zuckerman2023}, and $2821$ stars in archival HARPS spectra, where three
candidates in two M dwarfs and a red giant remain to be explained by the stars
themselves~\cite{FieldsGoodman2025}.

The spectroscopic surveys of the last decade observe many more nearby stars, at lower
resolution and over different wavelengths: the Gaia RVS spectra near the calcium triplet,
GALAH's four optical windows, APOGEE's $H$ band, which contains the $1530$ to $1565$\,nm
band in which Earth's fibre-optic lasers operate, and the wide optical coverage of DESI
and LAMOST. Their spectra are taken in many fibres at once and reduced in bulk, so a
search must tell a line in a star from the defects of fibres, detectors and sky
subtraction. This paper searches the nearby stars of all five surveys with methods that
compare each spectrum with others of the same survey.

\subsection{Contributions}
\begin{enumerate}[leftmargin=*,itemsep=1pt]
\item A search for unresolved emission lines in the spectra of nearby stars from Gaia
DR3, GALAH DR4, APOGEE DR17, DESI DR1 and LAMOST DR10, about $28000$ stars counted
survey by survey.
\item Comparison-spectrum models in which a narrow line that the comparison stars lack
cannot be fitted away, with each spectrum's own absorption-like residuals as the null.
\item The instrumental and terrestrial features of each survey that mimic a line, each
turned into a test.
\item The line flux at each survey's threshold, star by star, from the spectrum's noise
and the star's Gaia or 2MASS magnitude, as an isotropic luminosity and as the power of a
$10$\,m transmitter aimed at Earth.
\end{enumerate}

% =============================================================================
\section{Data and method}\label{sec:method}
% =============================================================================

\emph{The common method.} Each spectrum is divided by its upper envelope or continuum and
fitted, window by window, as a non-negative combination of the spectra of its most
similar stars in the same survey, with a signed continuum and signed first and second
derivatives of the fitted model, which absorb small differences of velocity and line
width. A narrow emission line in the target cannot be reproduced unless a comparison
spectrum carries it. The residual is searched with a matched filter of the instrument's
line profile, and peaks of both signs are kept: the absorption-like peaks arise from the
same imperfections of the model and the data as false emission-like ones, and serve as
their null. A peak of unresolved width is then
tested against the same position in other stars, in the stellar rest frame (a stellar
line the model misses) and in the observatory frame (sky, telluric and instrumental
features).

\emph{Gaia RVS.} The mean RVS spectra of Gaia DR3~\cite{Gaia2023,Katz2023}, $846$ to
$870$\,nm at a resolving power of $11500$, of the $2363$ FGK stars within $75$\,pc that
have one, are each fitted by their $25$ nearest neighbours in colour and absolute
magnitude. A positive peak counts only above every absorption-like peak of its
environment, and is tested against the cores of the \ion{Ca}{ii} triplet and the Paschen
lines, the same position in other stars, and deblended neighbours.

\emph{GALAH.} GALAH DR4~\cite{Buder2024} gives $699$ spectra of $621$ dwarfs within
$100$\,pc, at a resolving power of $28000$ in four windows from $471$ to $789$\,nm. Each
spectrum, camera by camera, is fitted by the $20$ sample spectra nearest in effective
temperature, gravity and metallicity; a candidate is then measured in the other fibres of
the same exposure.

\emph{APOGEE.} APOGEE DR17~\cite{Majewski2017,Abdurrouf2022} holds $H$-band spectra at a
resolving power of $22500$. For the $3370$ dwarfs within $100$\,pc with three or more
visits, the combined spectrum is fitted in $256$-pixel windows by its $12$ most similar
stars, and peaks under bright subtracted sky are OH airglow. Separately, each of the
$14178$ visits is divided by the star's median visit, which cancels anything present in
every visit, and lines that recur in several visits at one rest wavelength but at
different observatory wavelengths are grouped; detector dropouts that depress the median
are masked.

\emph{DESI.} The DESI DR1~\cite{DESI2025} Milky Way Survey~\cite{Cooper2023} spectra of
$7718$ stars with Gaia parallaxes of $10$ to $800$\,mas and S/N above $20$, $6266$ of
them cooler than $3500$\,K and $455$ white dwarfs, are fitted by their $12$ most similar
spectra; each candidate is checked exposure by exposure, since a coadd of several
exposures can carry a cosmic ray or a defect of one of them.

\emph{LAMOST.} LAMOST DR10's medium-resolution survey~\cite{Cui2012,Liu2020} observes
$495$ to $535$ and $630$ to $680$\,nm at a resolving power of $7500$ in several
exposures a night. In the $35179$ observations of $14400$ Gaia DR3 stars within
$100$\,pc with declination above $-15^\circ$ and $G<15$, a peak counts at $6\sigma$
above the local continuum with an unresolved width in a single exposure, and persists
when two exposures hold it within a third of a resolution element in the star's frame.
A feature is explained by an arm's end, by sharing with other stars in the stellar or
the instrument's frame, by recurrence in the same or neighbouring fibres of a
spectrograph, by instrument-frame recurrence on other nights, or by light in
neighbouring fibres of the same exposure, and what remains is set against the
coincidences of two exposures that chance makes and read again at its peak pixel in
every exposure of its star: a feature narrower than half the instrumental profile is a
detector pixel, one whose exposures peak more than half a resolution element apart is
not one feature, and one held by a single exposure of a night whose other exposures are
flat is a cosmic ray.

\emph{Flux limits.} Each search's threshold is a matched-filter significance $k$ in a
residual that is in units of the continuum: $k=5$ for the Gaia, APOGEE and DESI
searches, $5.5$ for GALAH and $6$ for a LAMOST exposure. An unresolved line whose
filter response sits at the threshold has a peak of $k\sigma/g$ of the continuum, with
$\sigma$ the spectrum's noise per pixel and $g$ the filter's gain, and a flux
$F=(k\sigma/g)\sqrt{2\pi}\,\sigma_\lambda F_\lambda$, with $\sigma_\lambda$ the
Gaussian width of the line profile and $F_\lambda$ the star's continuum flux density,
which we take from its Gaia BP, G and RP magnitudes, interpolated in wavelength in the
Vega system, and from its 2MASS $H$ magnitude for APOGEE. For Gaia the noise is that of
each star's residual; for GALAH, APOGEE and DESI it is the survey's per-pixel
signal-to-noise ratio of each camera or of the combined spectrum; for LAMOST it is the
per-exposure ratio of each arm, read from the headers of a random $1500$ of the
$35{,}179$ observations. An isotropic source at the star's Gaia distance $d$ then has
$L=4\pi d^2F$, and a transmitter aimed at Earth through a diffraction-limited $10$\,m
aperture needs $P=F d^2(\lambda/10\,\mathrm{m})^2$. These are thresholds, not
completeness limits: in the ultraviolet search of the same kind, lines injected into
archival spectra were recovered at $90$ per cent of the attainable fraction at $3$ to
$5$ times the nominal threshold, and the same factor applies here, since the same
envelopes, masks and comparison fits intervene.

\begin{table}[t]
\centering\small
\setlength{\tabcolsep}{4pt}
\caption{The surveys searched and their results. The null is the set of absorption-like
residual peaks found by the same method in the same spectra.}
\label{tab:surveys}
\begin{tabular}{l l r r >{\raggedright\arraybackslash}p{0.31\textwidth}}
\toprule
Survey & Range & $R$ & Stars & Result \\
\midrule
Gaia DR3 RVS & $846$--$870$\,nm & $11500$ & $2363$ & no line above the null \\
GALAH DR4 & $471$--$789$\,nm & $28000$ & $621$ & every candidate instrumental \\
APOGEE DR17, combined & $1.51$--$1.70$\um & $22500$ & $3370$ & no line above the null; one peak above it is OH airglow \\
APOGEE DR17, visits & $1.51$--$1.70$\um & $22500$ & $3370$ & $3431$ recurring emission-like groups against $3410$ absorption-like \\
DESI DR1 & $360$--$982$\,nm & $2000$--$5000$ & $7705$ & coadds of two or more exposures match the null \\
LAMOST DR10 MRS & $495$--$535$, $630$--$680$\,nm & $7500$ & $14400$ & $19$ two-exposure features, all detector pixels \\
\bottomrule
\end{tabular}
\end{table}

% =============================================================================
\section{Results}\label{sec:results}
% =============================================================================

Table~\ref{tab:surveys} summarises the results.

\emph{Gaia RVS and APOGEE.} No emission-like peak in the RVS spectra rises above the
absorption-like null. In the APOGEE combined spectra no emission-like peak exceeds the
null either in the continuum or inside stellar lines, once OH airglow is masked where the
visits' subtracted sky was bright; the one peak above the null was such airglow. In the
visit-to-visit search the recurring emission-like groups number $3431$ against $3410$
absorption-like ones, and the strongest, $48.2$, falls below the strongest absorption-like
group, $67.0$.

\emph{GALAH.} Every GALAH candidate is instrumental: pixels that several stars share in
the observatory frame, charge persistence in one fibre from a bright earlier exposure,
cosmic rays, and a ghost in the blue camera.

\emph{DESI.} The $1091$ stars with three or more exposures match the absorption-like null
in the continuum. Every peak above the null in them, and in the $1878$ coadds of two
exposures, vetted exposure by exposure, is a cosmic ray in one exposure, a corrupt
exposure, an OH residual, or a mismatch of the comparison spectra in a telluric band or
in a white dwarf's Paschen lines. The $4749$ stars with a single exposure, $62$ per cent
of the sample, have no second exposure to separate a line from a cosmic ray, so for them
the search ends at the peak list: their emission-like peaks are as numerous as their
absorption-like ones, and they are not searched in the sense the others are.

\emph{LAMOST.} The single exposures hold $12261$ unresolved spikes at $6\sigma$, and $110$
persist in two exposures in the star's frame. Arm ends, lines shared in the stellar or
the instrument's frame (sky residuals moved by the heliocentric correction), defects of
the same or neighbouring fibres, recurrence in the instrument's frame on other nights and
light in neighbouring fibres of the same exposures explain $91$; the $19$ left
(Table~\ref{tab:lamost}) lie in single stars, mostly in two of many exposures. The
strongest, at $241\sigma$ in three exposures of one night of a K dwarf, lies a few
pixels from a block of pixels at $1800$ times the continuum in all five exposures of
that night, with the pipeline's variance there inflated a thousandfold: a defect of
those extractions, absent from the star's $178$ other exposures, against which the
residual's scatter was measured. The other $18$ are as many as the $17$ coincidences
of two exposures that chance makes, and read at the peak pixel in every exposure of
their stars, all $19$ are detector pixels: each rises above half its excess over one or
two pixels, where the instrumental profile's full width at half maximum is $5.6$ to
$6.5$ pixels, so no light passing through the spectrograph made it. Six of them sit as
well on masked, variance-inflated or dead pixels, eight have exposures whose peaks lie
more than half a resolution element apart, four fill one night and none of the star's
other nights, and three are single exposures of nights whose other exposures are flat.

\emph{Flux limits.} Table~\ref{tab:limits} gives the line flux at each survey's
threshold. The near-infrared and the deep DESI coadds reach $2\times10^{-16}$\,erg\,cm$^{-2}$\,s$^{-1}$
at the median star and the optical surveys of brighter stars $1$ to
$3\times10^{-15}$; as isotropic luminosities at the stars' distances these are $10^{19}$
to $2\times10^{20}$\,W, and as transmitters behind a $10$\,m aperture aimed at Earth
$5$ to $70$\,kW at the median star, $300$\,W at the DESI stars within $20$\,pc and $5$
to $50$\,W at the nearest star of each survey. The thresholds of Tellis and Marcy's
search, megawatts for a $10$\,m aperture at their stars' distances~\cite{TellisMarcy2017},
are reached here at every survey's median star, with the three-to-five-fold completeness
factor of the previous paragraph still to be applied.

\begin{table}[tp]
\centering\small
\setlength{\tabcolsep}{4pt}
\caption{The line flux at each search's threshold: the median over the stars (for
GALAH, DESI and LAMOST over the cameras, arms and exposures) of $F$, of the isotropic
luminosity $L$ at the star's distance and of the power $P_{10\,\rm m}$ of a $10$\,m
transmitter aimed at Earth; the number of stars within $20$\,pc; and $P_{10\,\rm m}$ at
the nearest star. The LAMOST row is from the $1163$ stars among a random $1500$
observations whose headers were read. APOGEE's $P_{10\,\rm m}$ is larger than DESI's at
a smaller flux because the beam at $1.6$\um\ is five times wider.}
\label{tab:limits}
\begin{tabular}{l r r r r r r}
\toprule
Survey & Stars & $F$ (erg\,cm$^{-2}$\,s$^{-1}$) & $L$ (W) & $P_{10\,\rm m}$ (W) & $N_{<20\,\rm pc}$ & $P_{10\,\rm m}$, nearest (W) \\
\midrule
Gaia DR3 RVS & $2363$ & $2.9\times10^{-15}$ & $1.2\times10^{20}$ & $7\times10^{4}$ & $27$ & $3.5\times10^{4}$ ($7.2$\,pc) \\
GALAH DR4 & $621$ & $1.3\times10^{-15}$ & $1.1\times10^{20}$ & $3\times10^{4}$ & $4$ & $50$ ($4.3$\,pc) \\
APOGEE DR17 & $2931$ & $1.6\times10^{-16}$ & $1.1\times10^{19}$ & $2\times10^{4}$ & $4$ & $2.5\times10^{3}$ ($13.8$\,pc) \\
DESI DR1 & $7718$ & $2.1\times10^{-16}$ & $1.4\times10^{19}$ & $5\times10^{3}$ & $96$ & $5$ ($3.6$\,pc) \\
LAMOST DR10 MRS & $1163$ & $3.0\times10^{-15}$ & $2.0\times10^{20}$ & $5\times10^{4}$ & $20$ & $4\times10^{3}$ ($16$\,pc) \\
\bottomrule
\end{tabular}
\end{table}

\begin{table}[tp]
\centering\small
\setlength{\tabcolsep}{5pt}
\caption{The $19$ LAMOST features that pass the tests against other stars, fibres and
nights: the star, the arm, the rest and observed wavelengths (air), the exposures
holding the feature against the star's exposures, the nights they span, and the largest
single-exposure significance. Each is one or two pixels wide.}
\label{tab:lamost}
\begin{tabular}{l l r r r r r}
\toprule
Gaia DR3 & Arm & $\lambda_{\rm rest}$ (\AA) & $\lambda_{\rm obs}$ (\AA) & Exposures & Nights & $z$ \\
\midrule
4420231361513433344 & B & 4947.67 & 4947.83 & 2/6 & 1 & 9.5 \\
1407678284339274368 & B & 5102.09 & 5102.55 & 2/8 & 1 & 9.4 \\
1407678284339274368 & B & 5134.99 & 5135.46 & 2/8 & 1 & 8.3 \\
608182831000585728 & B & 5269.81 & 5270.60 & 2/30 & 1 & 7.8 \\
815877833000303488 & B & 5290.17 & 5289.98 & 2/238 & 2 & 9.1 \\
1295249440828061696 & B & 5331.15 & 5331.32 & 3/160 & 1 & 20.0 \\
1894835705805757824 & B & 5344.89 & 5345.21 & 2/152 & 2 & 17.8 \\
64953069592851328 & R & 6307.36 & 6306.18 & 2/124 & 2 & 11.1 \\
1898531439263850240 & R & 6322.92 & 6323.26 & 5/106 & 1 & 14.5 \\
814125692502763904 & R & 6335.08 & 6335.53 & 2/214 & 2 & 22.4 \\
687208343753227520 & R & 6433.52 & 6433.98 & 2/52 & 2 & 28.2 \\
1547717514092903424 & R & 6551.41 & 6550.29 & 2/196 & 2 & 9.9 \\
2556424771496343424 & R & 6626.80 & 6626.62 & 2/76 & 2 & 9.0 \\
1915323176544379520 & R & 6728.65 & 6729.16 & 2/180 & 2 & 11.7 \\
779262652684989312 & R & 6756.41 & 6756.55 & 2/174 & 2 & 6.5 \\
3872246689058776832 & R & 6764.64 & 6774.04 & 2/123 & 2 & 33.0 \\
1386361674454950912 & R & 6772.31 & 6772.58 & 3/183 & 1 & 241.4 \\
161252970877957632 & R & 6783.43 & 6783.15 & 4/34 & 1 & 7.0 \\
986071444021690112 & R & 6801.65 & 6802.17 & 4/22 & 1 & 21.1 \\
\bottomrule
\end{tabular}
\end{table}

% =============================================================================
\section{Related work}\label{sec:related}
% =============================================================================

\emph{Laser searches.} The searches of Section~\ref{sec:intro} use high-resolution
spectra of radial-velocity programmes~\cite{ReinesMarcy2002,TellisMarcy2015,TellisMarcy2017,Zuckerman2023,FieldsGoodman2025},
and the HARPS search, the largest of them, ends as this one does with a handful of
features in stars whose own spectra are the first thing to exclude. We find no earlier
laser search of Gaia, GALAH, APOGEE, DESI or LAMOST spectra.

% =============================================================================
\section{Conclusion}\label{sec:conclusion}
% =============================================================================

No unresolved emission line that a laser aimed at Earth could produce is found in the
nearby stars of five large spectroscopic surveys, beyond what the surveys' instruments,
the sky and chance produce. The surveys extend the earlier laser searches to tens of
thousands of stars, to the near infrared and to the telecommunications band.

The limits are thresholds in each spectrum's noise, turned into fluxes by the stars'
magnitudes, and apply to lines that are steady in the combined spectra or that recur in
the exposures available: $10^{19}$ to $10^{20}$\,W isotropic at the median star,
kilowatts to tens of kilowatts behind a $10$\,m aperture, watts at the nearest stars.
A line inside a strong stellar line, one in the telluric bands that
the models fit away, or one that was off when the spectra were taken is not constrained;
the $62$ per cent of DESI stars with a single exposure need a second, and injecting
lines into the archival spectra, as the ultraviolet search did, would turn these
thresholds into completeness limits survey by survey.

\section*{Acknowledgements}

This work uses data from the European Space Agency mission Gaia, processed by the Gaia
Data Processing and Analysis Consortium; from the GALAH survey, based on observations at
the Anglo-Australian Telescope; from SDSS-IV's APOGEE; from the Dark Energy Spectroscopic
Instrument, retrieved through NOIRLab's SPARCL; and from the Large Sky Area Multi-Object
Fiber Spectroscopic Telescope (LAMOST), funded by the National Development and Reform
Commission and operated by the National Astronomical Observatories, Chinese Academy of
Sciences.

\begin{thebibliography}{99}\small\raggedright
\setlength{\itemsep}{1pt}

\bibitem{SchwartzTownes1961}
R.~N. Schwartz and C.~H. Townes, Interstellar and interplanetary communication by
optical masers, \emph{Nature} \textbf{190}, 205 (1961).

\bibitem{ReinesMarcy2002}
A.~E. Reines and G.~W. Marcy, Optical search for extraterrestrial intelligence: a
spectroscopic search for laser emission from nearby stars, \emph{Publ. Astron. Soc.
Pac.} \textbf{114}, 416 (2002).

\bibitem{TellisMarcy2015}
N.~K. Tellis and G.~W. Marcy, A search for optical laser emission using Keck HIRES,
\emph{Publ. Astron. Soc. Pac.} \textbf{127}, 540 (2015).

\bibitem{TellisMarcy2017}
N.~K. Tellis and G.~W. Marcy, A search for laser emission with megawatt thresholds from
5600 FGKM stars, \emph{Astron. J.} \textbf{153}, 251 (2017).

\bibitem{Zuckerman2023}
A.~Zuckerman, Z.~Ko, H.~Isaacson, et al., The Breakthrough Listen search for
intelligent life: a laser search pipeline for the Automated Planet Finder,
\emph{Astron. J.} \textbf{165}, 114 (2023).

\bibitem{FieldsGoodman2025}
B.~Fields and J.~C. Goodman, A 2821 star optical SETI survey using ESO HARPS archival
data, arXiv:2508.08628 (2025).

\bibitem{Gaia2023}
Gaia Collaboration, A.~Vallenari, A.~G.~A. Brown, et al., Gaia Data Release 3: summary
of the content and survey properties, \emph{Astron. Astrophys.} \textbf{674}, A1 (2023).

\bibitem{Katz2023}
D.~Katz, P.~Sartoretti, A.~Guerrier, et al., Gaia Data Release 3: properties and
validation of the radial velocities, \emph{Astron. Astrophys.} \textbf{674}, A5 (2023).

\bibitem{Buder2024}
S.~Buder, J.~Kos, E.~X. Wang, et al., The GALAH survey: data release 4,
arXiv:2409.19858 (2024).

\bibitem{Majewski2017}
S.~R. Majewski, R.~P. Schiavon, P.~M. Frinchaboy, et al., The Apache Point Observatory
Galactic Evolution Experiment (APOGEE), \emph{Astron. J.} \textbf{154}, 94 (2017).

\bibitem{Abdurrouf2022}
Abdurro'uf, K.~Accetta, C.~Aerts, et al., The seventeenth data release of the Sloan
Digital Sky Surveys, \emph{Astrophys. J. Suppl.} \textbf{259}, 35 (2022).

\bibitem{DESI2025}
DESI Collaboration, M.~Abdul-Karim, A.~G. Adame, et al., Data release 1 of the Dark
Energy Spectroscopic Instrument, arXiv:2503.14745 (2025).

\bibitem{Cooper2023}
A.~P. Cooper, S.~E. Koposov, C.~Allende Prieto, et al., Overview of the DESI Milky Way
Survey, \emph{Astrophys. J.} \textbf{947}, 37 (2023).

\bibitem{Cui2012}
X.-Q. Cui, Y.-H. Zhao, Y.-Q. Chu, et al., The Large Sky Area Multi-Object Fiber
Spectroscopic Telescope (LAMOST), \emph{Res. Astron. Astrophys.} \textbf{12}, 1197
(2012).

\bibitem{Liu2020}
C.~Liu, J.~Fu, J.~Shi, et al., LAMOST medium-resolution spectroscopic survey
(LAMOST-MRS): scientific goals and survey plan, arXiv:2005.07210 (2020).

\end{thebibliography}

\end{document}
